\documentclass[12pt,a4paper]{article}
\usepackage[utf8]{inputenc}
\usepackage[vietnamese]{babel}
\usepackage[margin=2cm]{geometry}
\usepackage{tcolorbox}
\usepackage{enumitem}
\usepackage{booktabs}
\usepackage{fancyhdr}

\pagestyle{fancy}
\fancyhf{}
\lhead{\textbf{CÔNG AN THÀNH PHỐ - PHÒNG CSHS}}
\rhead{\textbf{CHUYÊN ÁN TEST-99}}
\cfoot{Trang \thepage}

\begin{document}

\begin{center}
    \textbf{\Large CỘNG HÒA XÃ HỘI CHỦ NGHĨA VIỆT NAM}\\
    \textbf{\large Độc lập - Tự do - Hạnh phúc}\\
    \rule{8cm}{0.5pt}\\[0.5cm]
    
    \textbf{\LARGE BÁO CÁO KHÁM NGHIỆM TỬ THI SƠ BỘ}\\
    \textbf{Số hiệu:} \texttt{BKNTT-2026/0725-01}
\end{center}

\vspace{0.3cm}

\begin{tcolorbox}[colback=gray!10!white,colframe=black!75!black,title=\textbf{THÔNG TIN HÀNH CHÍNH}]
\begin{itemize}[leftmargin=*]
    \item \textbf{Cơ quan thực hiện:} Phòng Kỹ thuật Hình sự - Công an Thành phố
    \item \textbf{Thời gian khám nghiệm:} 09:30 ngày 25/07/2026
    \item \textbf{Địa điểm thực hiện:} Hiện trường nhà số 14, Đường Bờ Sông, Phân khu Cảng phía Bắc
\end{itemize}
\end{tcolorbox}

\section*{I. THÔNG TIN NẠN NHÂN}
\begin{itemize}
    \item \textbf{Họ và tên:} Nguyễn Văn Khang
    \item \textbf{Năm sinh:} 1991 (35 tuổi)
    \item \textbf{Giới tính:} Nam
    \item \textbf{Tình trạng nhận dạng:} Thi thể nguyên vẹn, đang trong giai đoạn bắt đầu phân hủy nhẹ.
\end{itemize}

\section*{II. KẾT QUẢ KHÁM NGHIỆM BÊN NGOÀI}

\subsection*{1. Vết thương vùng Đầu / Vai / Lưng}
\begin{itemize}
    \item Có 01 vết bầm tím lớn kích thước \texttt{6cm x 4cm} vùng sau chẩm (gáy).
    \item Vùng vai phải có vết bầm trợt da nhẹ do va đập vào bề mặt cứng/góc cạnh.
    \item \textbf{Đánh giá sơ bộ:} Tổn thương do tác động lực tụt ngã/xô ngã đập vào bề mặt gỗ hoặc mép tủ. Vết thương này gây chấn thương sọ não nhẹ và \textbf{ngất/bất tỉnh ngắn}, không phải nguyên nhân trực tiếp gây tử vong.
\end{itemize}

\subsection*{2. Vết thương vùng Cổ / Ngực (Tổn thương chết người)}
\begin{itemize}
    \item Có 01 vết đâm rách sâu vùng cổ bên trái kích thước \texttt{3.5cm}, bờ vết thương sắc gọn, gây đứt động mạch cảnh ngoài.
    \item Lượng máu chảy ra hiện trường xung quanh thi thể ước tính lớn (từ 1.2L đến 1.5L).
    \item \textbf{Đánh giá sơ bộ:} Nạn nhân tử vong do \textbf{mất máu cấp (sốc mất máu)} do vật sắc nhọn đâm trực tiếp.
\end{itemize}

\section*{III. THỜI GIAN TỬ VONG ƯỚC TÍNH}
Dựa trên độ cứng tử thi (\textit{rigor mortis}) và nhiệt độ môi trường căn nhà đóng kín:
\begin{itemize}
    \item \textbf{Khung giờ tử vong ước tính:} Từ \textbf{20:30 đến 21:30 ngày 24/07/2026}.
\end{itemize}

\section*{IV. KẾT LUẬN SƠ BỘ}
\begin{enumerate}
    \item Nguyên nhân tử vong trực tiếp: Sốc mất máu cấp do vết đâm sắc nhọn vùng cổ đứt động mạch.
    \item Nạn nhân có dấu hiệu giằng co/xô ngã trước đó (gây vết bầm vùng gáy).
    \item Hướng vết đâm có góc từ trên xuống dưới, thực hiện khi nạn nhân đang trong tư thế nằm ngửa hoặc gục trên sàn nhà.
\end{enumerate}

\vspace{1cm}

\begin{tcolorbox}[colback=red!5!white,colframe=red!75!black,title=\textbf{XÁC NHẬN BÁO CÁO}]
\begin{center}
    \textit{Bác sĩ Pháp y Giám định chính}\\
    \vspace{1.5cm}
    \textbf{BS. Lê Văn Hoàng} (Đã ký và đóng dấu)
\end{center}
\end{tcolorbox}

\end{document}
